\begin{proofbox}
	\textbf{\textcolor{CProof}{思路提示}}

	利用对称变换，将同侧两定点问题转化为异侧，再依据两点间线段最短得证。

	\medskip
	\textbf{\textcolor{CProof}{正式推导}}
	\begin{description}[style=nextline, leftmargin=2.8em, labelsep=0.5em, itemsep=0.55em]
		\item[\textbf{步骤一（作对称点）：}] 作$A$关于直线$l$的对称点$A'$，设$l$交$AA'$于点$O$，则$AO=OA'$且$AA'\perp l$。
			\[
				AO = OA',\quad AA'\perp l.
			\]
			\par\textbf{理由：} 对称变换保证$l$垂直平分$AA'$。

		\item[\textbf{步骤二（转化线段和）：}] 对于$l$上任意一点$Q$，有$QA=QA'$，从而
			\[
				QA+QB = QA'+QB.
			\]
			\par\textbf{理由：} 等量代换。

		\item[\textbf{步骤三（利用线段公理）：}] 连接$A'B$，由两点间线段最短得
			\[
				QA'+QB \ge A'B.
			\]
			\par\textbf{理由：} 两点之间线段最短。

		\item[\textbf{步骤四（取等条件）：}] 当$Q$在线段$A'B$上时，等号成立。记$A'B$与$l$的交点为$P$，则当$Q=P$时，$PA+PB$取最小值$A'B$。
			\[
				(PA+PB)_{\min}=A'B.
			\]
			\par\textbf{理由：} 三点共线时折线化为直线段。
	\end{description}

	\medskip
	\textbf{\textcolor{CProof}{结论回扣}}

	\textbf{将军饮马模型的核心是“对称转化”，最终将$PA+PB$的最小值转化为$A'B$的长度。}
\end{proofbox}
